\section{Policy Network Architectures}

Financial markets exhibit strong temporal dependencies, non-stationary behavior, and 
long-range interactions across assets and indicators. Consequently, the policy network
must not only learn the mapping from states to actions, but also capture temporal market
structure.

To investigate the effect of temporal modeling in DRL trading, this work evaluates four
policy families: standard MLP policies, recurrent LSTM policies, and two custom 
Transformer-based feature extractors integrated into the Stable-Baselines3 framework.

Standard MLP policies treat each observation independently:
\[
a_t = \pi(s_t)
\]
thus ignoring historical information. While computationally efficient, this assumption
is restrictive in financial environments where decisions depend on previous market states.

LSTM-based policies extend this formulation by introducing a hidden state:
\[
h_t = f(h_{t-1}, s_t)
\]
allowing information propagation through time. However, long temporal dependencies must
be compressed into a fixed-size memory representation, which may lead to information loss.

Transformer architectures overcome this limitation through self-attention, allowing each
timestep to directly interact with all previous observations inside a lookback window.
This enables explicit modeling of delayed market reactions, regime changes and cross-asset
dependencies, making Transformers particularly suitable for multi-asset trading tasks.

\subsection{Transformer Feature Extractor (Full Fusion Architecture)}

The first proposed architecture is a fully Transformer-based feature extractor where
portfolio and market information are processed jointly.

For each timestep $t$, the observation consists of the concatenation of portfolio state
and market features:
\[
x_t =
[s_t^{portfolio}, s_t^{market}]
\]

A temporal window of length $T$ is constructed:
\[
X =
[x_{t-T+1}, \dots, x_t]
\]
which forms the Transformer input.

Each observation is projected into a latent representation of dimension
$d_{model}=128$:
\[
z_t = W_{proj}x_t
\]
followed by addition of learnable positional embeddings:
\[
\tilde{z}_t = z_t + p_t
\]
to preserve temporal ordering.

The resulting sequence is processed using a Transformer encoder composed of two encoder
layers with four attention heads, feedforward dimension $256$, pre-layer normalization,
and dropout $0.1$.

The final latent representation is obtained from the output corresponding to the most
recent timestep:
\[
h_t = Transformer(X)_T
\]
followed by Layer Normalization.

This representation $h_t$ is used as a compact feature vector and is passed to the
actor and critic networks of the underlying DRL algorithm to produce the policy
distribution and value estimates.

This architecture performs full feature fusion, allowing attention to jointly model
portfolio evolution, market dynamics and cross-asset interactions within a unified
representation space.

\subsection{Hybrid MLP--Transformer Architecture}

The second architecture introduces an explicit separation between sequential market
information and instantaneous portfolio state.

The market branch receives only market observations over the temporal window:
\[
X^{market}
=
[s_{t-T+1}^{market}, \dots, s_t^{market}]
\]
which are projected to $d_{model}=128$, combined with learnable positional embeddings
and processed by the same Transformer encoder configuration used in the previous model.

The resulting market embedding is extracted from the final timestep output:
\[
h_t^{market}
=
Transformer(X^{market})_T
\]

In parallel, portfolio information is processed independently using an MLP branch.
The input consists only of the most recent portfolio state:
\[
h_t^{portfolio}
=
MLP(s_t^{portfolio})
\]
implemented using two fully connected layers of size $64$ with ReLU activations.

The two latent representations are fused through concatenation:
\[
z_t
=
[h_t^{market}, h_t^{portfolio}]
\]
followed by a fully connected fusion layer projecting the representation back to
$d_{model}=128$.

The resulting fused embedding $z_t$ is passed to the actor and critic networks of the
reinforcement learning algorithm, where it is used to compute action distributions and
value estimates.

This design is motivated by the different characteristics of the two information sources.
Market observations are inherently sequential and benefit from temporal attention
mechanisms, whereas portfolio information represents an instantaneous allocation state
that does not require sequence modeling. Separating the branches reduces unnecessary
coupling while preserving expressive capacity through the fusion stage.